{
\def\AMCbeginQuestion#1#2{\par\noindent{\bf (POSCOMP 2023 - Questão 04)}}
\begin{question}{} Nos jogos da Mega-Sena, são sorteados a cada concurso 6 números no intervalo de 1 a 60. Leva o prêmio quem acertar os 6 números sorteados. O apostador, ao fazer um jogo, pode optar por preencher um bilhete com 6, 7, 8 ou 9 números escolhidos. Se o apostador registra um bilhete com 8 números escolhidos, quantos bilhetes de 6 números ele faria com os mesmos 8 números escolhidos?
\begin{choices}
\correctchoice{28}
\wrongchoice{56}
\wrongchoice{$8!$}
\wrongchoice{$6!$}
\wrongchoice{72}
\end{choices}
\end{question}
}
%
{
\def\AMCbeginQuestion#1#2{\par\noindent{\bf (POSCOMP 2023 - Questão 05)}}
\begin{question}{} Recentemente, com a pandemia de Covid-19, houve grande interesse em determinar conjuntos de regiões (de países, estados, municípios, etc.) com alta incidência da doença, com o objetivo de determinar políticas de mitigação da doença nesses locais. Nesse sentido, dado um mapa subdividido em regiões, um cluster é definido como sendo um subconjunto de regiões desse mapa (nesse caso, pode ser formado por regiões que não fazem fronteira entre si). Qual o número de possíveis clusters para um mapa com 10 regiões?
\begin{choices}
\correctchoice{1024}
\wrongchoice{10}
\wrongchoice{100}
\wrongchoice{512}
\wrongchoice{20}
\end{choices}
\end{question}
}
%
{
\def\AMCbeginQuestion#1#2{\par\noindent{\bf (POSCOMP 2023 - Questão 06)}}
\begin{question}{} Um anagrama é uma nova palavra formada pela permutação de letras de uma palavra. Essa nova palavra não precisa fazer sentido. Um anagrama de POSCOMP, por exemplo, seria MOCPSOP. Quantos são os anagramas da palavra POSCOMP?
\begin{choices}
\correctchoice{1260}
\wrongchoice{126}
\wrongchoice{252}
\wrongchoice{2520}
\wrongchoice{5040}
\end{choices}
\end{question}
}
